% Source: https://artofproblemsolving.com/community/c3h1252036_distance_to_triangle_sides_from_bisectors_bases
% Stable identity: aops-topic-1252036-q01
For a point $X$ of the closed triangle, let $d_{AB}(X)$, $d_{AC}(X)$, $d_{BC}(X)$ be its distances to the three side lines. Inside the triangle each distance is an affine (linear plus constant) function of the coordinates of $X$, because the triangle lies on one side of each line.

\begin{asy}
usepackage("amsmath");
size(4.5cm);
defaultpen(fontsize(10pt));
pair A=(0,0),B=(6,0),C=(2,4);
real a=abs(B-C),bb=abs(A-C),c=abs(A-B);
pair A1=(bb*B+c*C)/(bb+c), B1=(a*A+c*C)/(a+c);
draw(A--B--C--cycle,linewidth(1)); draw(A--A1,gray); draw(B--B1,gray); draw(A1--B1,blue);
pair M=0.4*A1+0.6*B1; dot(M);
pair F1=A+dot(M-A,unit(B-A))*unit(B-A), F2=A+dot(M-A,unit(C-A))*unit(C-A), F3=B+dot(M-B,unit(C-B))*unit(C-B);
draw(M--F1,red+dashed); draw(M--F2,red+dashed); draw(M--F3,red+dashed);
label("$A$",A,SW);label("$B$",B,SE);label("$C$",C,N);label("$A_1$",A1,NE);label("$B_1$",B1,W);label("$M$",M,S);
label("\mbox{}",(min(currentpicture,user=true)+max(currentpicture,user=true))/2,black); // keeps the text after the figure black when LaTeX typesets the labels inline
\end{asy}

Consider $g(X)=d_{AB}(X)-d_{AC}(X)-d_{BC}(X)$, also affine. At $A_1$: $A_1$ lies on $BC$, so $d_{BC}(A_1)=0$, and on the bisector from $A$, so $d_{AB}(A_1)=d_{AC}(A_1)$; thus $g(A_1)=0$. At $B_1$: $d_{AC}(B_1)=0$ and $d_{AB}(B_1)=d_{BC}(B_1)$; thus $g(B_1)=0$. An affine function that vanishes at both ends of a segment vanishes on the whole segment (along the segment it is a linear function of one parameter). Hence $g(M)=0$ for all $M$ on $A_1B_1$, which is the claim. $\blacksquare$
